\section{Directed rounding without activation margins}
\label{sec:revisit-rounding}
This section changes the exact recurrence to a specified rounded algorithm.
The scalar rebase idea is adapted from the active manuscript's bounded
arithmetic appendix. The new features are one-sided neighbor records,
an orthant projection, a coarser kinetic grid, and a rounded cooling
schedule. All assertions concern the actual stored states.

Let $h$ be a reciprocal power of two and let
$\lfloor a\rfloor_h=h\lfloor a/h\rfloor$. Retain the parameters
$\theta,\mu,\chi,\eta$ above. Write
$c=(1-\alpha)/2$ and $q_0=(1+\alpha)/2$.
At iteration boundaries store
\[
 x_i/\omega_i=\sigma X_i,\qquad Z_i=z_i/\omega_i,
 \qquad \sigma,X_i\in h\sN_0,\quad
 Z_i\in(h/\theta)\sN_0,\quad 1/2\leq\sigma\leq1.
\]
The kinetic grid is deliberately $h/\theta$: its smallest positive
coordinate contributes at least one primal grid unit before scaling.
Store neighbor records $L_i\in h\sN_0$ with the invariant
\begin{equation}
 0\leq\sum_{j\sim i}X_j-L_i\leq4d_i h.
 \label{eq:revisit-directed-neighbors}
\end{equation}
Kinetic neighbor sums are exact. All degrees in these formulas are original
graph degrees, including at the boundary.

\paragraph{One rounded step.}
In the reporter replace the normalized primal response by
$(q_0-\mu)X_i-cL_i/d_i$ and leave all other terms exact. If $\vq$ is
the ideal raw vector at the actual stored states, the represented raw
vector is $\vq+\vu$ where
\begin{equation}
 -\frac{2h}{\theta}\vomega\leq\vu\leq\vzero.
 \label{eq:revisit-directed-raw}
\end{equation}
Let $\vp_0=[\vq+\vu]_+$ be an analytical comparison vector and set
\[
 p_i/\omega_i=\left\lfloor p_{0,i}/\omega_i\right\rfloor_{h/\theta}.
\]
The reporter enumerates the closed tail
$k_i\geq(\lambda/\theta+h/\theta)/\sigma$. It therefore opens exactly
the positive stored kinetic coordinates, without enumerating smaller
positive ideal coordinates.

Next set $s=\lfloor\chi\sigma\rfloor_h$.
If $s<1/2$, first rebase the \emph{old} primal records by
\begin{equation}
 X_i\gets\lfloor sX_i\rfloor_h,\qquad
 L_i\gets\max\{0,\lfloor sL_i\rfloor_h-d_i h\},
 \qquad s\gets1.
 \label{eq:revisit-directed-rebase}
\end{equation}
These replacements use the old values and make no row scans. Then, on
the emitted coordinates only, perform
\begin{equation}
 X_i\gets X_i+\left\lfloor\frac{\theta(p_i/\omega_i)}s\right\rfloor_h,
 \qquad \sigma\gets s,\qquad \vz\gets\vp.
 \label{eq:revisit-rounded-primal}
\end{equation}
Scatter the actual increments in $X_i$ and form exact kinetic neighbor
sums from these same charged rows. Rebuild all keys after a rebase;
otherwise refresh only the usual touched coordinates and exceptions.
No query occurs while an update is incomplete.
An old positive primal coordinate can round to zero at a rebase. Thus
actual rounded supports need not be nested; the implementation retains a
nested history of exposed labels, which is what the storage and rebase
charges use.

\begin{lemma}[Directed records and kinetic mass]
\label{lem:revisit-rounded-records}
Suppose $h\leq1/16$ and the incoming primal mass is at most three.
The updates preserve \cref{eq:revisit-directed-neighbors} and satisfy
\begin{equation}
 \vp=\vp_0+\ve_p,\quad
 \vx^+=\chi\vx+\theta\vp+\ve_x,\qquad
 -\frac h\theta\vomega\leq\ve_p\leq0,
 \quad -8h\vomega\leq\ve_x\leq0.
 \label{eq:revisit-rounded-errors}
\end{equation}
Also $\theta\vomega^\trans\vp\leq2\vomega^\trans\vx^+$.
\end{lemma}
\begin{proof}
At a rebase let the scalar floor losses be $\ell_j,\ell_i^L\in[0,h)$.
Before the nonnegative clamp, the new neighbor deficit equals
\[
 s\Big(\sum_{j\sim i}X_j-L_i\Big)
                -\sum_{j\sim i}\ell_j+\ell_i^L+d_i h.
\]
It is nonnegative and at most $4s d_i h+2d_i h\leq4d_i h$.
The clamp can only reduce that nonnegative deficit. Actual scattered
increments preserve it between rebases. Substitution into the reporter
gives \cref{eq:revisit-directed-raw}.

The incoming density bound gives $X_i\leq6$. Rounding the scale loses
at most $6h$ in physical primal density. A rebase loses at most another
$h$, and the emitted-coordinate increment floor loses at most $h$.
All losses are downward, proving \cref{eq:revisit-rounded-errors}.
The scale is positive because $\chi\sigma\geq1/4$.
If $p_i>0$, then $\theta p_i/\omega_i\geq h$ and $s\leq1$.
The floor in \cref{eq:revisit-rounded-primal} retains at least half its
argument. The old primal contribution is nonnegative, so
$x_i^+\geq\theta p_i/2$ coordinatewise. The rebase precedes this
increment and cannot erase it.
\end{proof}

\begin{lemma}[Two perturbed energies without a mass cap]
\label{lem:revisit-rounded-energies}
At an incoming state with
$\vomega^\trans\vx\leq3$ and
$\vomega^\trans\vz\leq6/\theta$, either comparison energy obeys
\begin{equation}
 E^+\leq\chi E+256h,
 \qquad B_{\rm fixed}^+\leq\chi B+256h.
 \label{eq:revisit-rounded-energy-step}
\end{equation}
The first comparator is the current RPPR optimum, and the second is
$\vt=\mQ^{-1}\vb$ as before.
\end{lemma}
\begin{proof}
Both comparison vectors have weighted mass at most one.
The identity
\[
 \vq=\frac{\mI-\mQ}{1+\theta}\vz
       -\frac{\mQ-\mu\mI}{\theta(1+\theta)}\vx
       +\frac{\vb-\lambda\vomega}{\theta}
\]
and $\vu\leq0$ imply
\[
 \vomega^\trans\vp_0
 \leq\frac{1-\alpha}{1+\theta}\vomega^\trans\vz
       +\frac{c}{\theta(1+\theta)}\vomega^\trans\vx
       +\frac\alpha\theta\leq9/\theta.
\]
Here the diagonal of $\mQ-\mu\mI$ is nonnegative. Thus the ideal
next primal mass is at most $12$, and the absolute weighted difference
between $\vp_0$ and either comparator is at most $10/\theta$.

In the imported comparison identity, replacing the projection normal of
$\vq+\vu$ by that of $\vq$ costs at most
$\mu(10/\theta)(2h/\theta)=20h$. This estimate also holds in the
$\mQ$ metric because $|\mQ|\vomega=\vomega$.
Rounding $\vp_0$ downward costs at most $10h+9h/2$ in either mirror
term: use the coordinate error $h/\theta$, the mass bound $9/\theta$,
and $\mu=\theta^2$.

The total primal loss relative to $\chi\vx+\theta\vp_0$ has density
at most $9h$ and mass at most $12$. Its squared norm is at most
$12(9h)$. The absolute weighted mass of the residual at that ideal
primal point is at most $13$. Expanding either quadratic therefore costs
at most $(13+6)(9h)=171h$. The positive linear mass penalties decrease
under the downward loss. The three costs sum to less than $256h$.
\end{proof}

\paragraph{Finite cooling parameters.}
For rational $\rho$, set $g=\theta\rho/8$. Start at $r_0=1/d_v$ and use
\begin{equation}
 r_{k+1}=\max\{\rho,g\lceil\eta r_k/g\rceil\}.
 \label{eq:revisit-rounded-cooling}
\end{equation}
Since $\rho/g=8/\theta$ is an integer, all parameters after the first
are integer multiples of the same $g$. Whenever $r_k>\rho$,
\[
 \eta r_k\leq r_{k+1}
 \leq\max\{\rho,(1-3\theta/8)r_k\}<r_k.
\]
Thus the finite ramp terminates and
$\sum_{\rm ramp}1/r_k<3/(\theta\rho)$.
There is no growing exact geometric-power denominator.

\begin{theorem}[Rounded two-phase semantic completion]
\label{thm:revisit-rounded-semantic}
Put $\rho=\epsppr/2$, $\tau=\alpha\epsppr^2/32$, and choose the
coarsest dyadic grid satisfying
\begin{equation}
 h\leq\min\{1/16,\theta,
           \theta\alpha^2\rho/256,\theta\tau/512\}.
 \label{eq:revisit-rounded-grid}
\end{equation}
Use the rounded operations above, with the same constant-output branches
as \cref{alg:revisit-ppr}. After the ramp, let $q$ be the least nonnegative
integer with $2^q\geq416/\epsppr$ and run $N=q/\theta$ frozen rounded
steps. Determine $q$ by integer doubling; no growing contraction power
is stored. Final exact clipping of the represented
densities by $\epsppr/4$ satisfies \cref{eq:revisit-semantic}.
With $N$ frozen steps, cumulative kinetic volume is at most
\begin{equation}
 60\sum_{k<K}1/r_k
 \leq\frac{60}{\rho}(N+3/\theta),
 \qquad N\leq\theta^{-1}\{1+\log_2(416/\epsppr)\}.
 \label{eq:revisit-rounded-work}
\end{equation}
\end{theorem}
\begin{proof}
Set $\Gamma=256h/\theta$. The rounded bank induction is
\[
 B_{k+1}\leq\chi B_k+256h+\alpha(\lambda_k-\lambda_{k+1}),
 \qquad B_k\leq\alpha\lambda_k+\Gamma.
\]
The latter follows from $\chi\Gamma+256h=\Gamma$ and the same
cooling-ratio inequality as in \cref{thm:revisit-cooling-bank}.
Since $\Gamma\leq\alpha^2\rho\leq\alpha\lambda_k$, the new primal
mass is at most three. \Cref{lem:revisit-rounded-records} then bounds
the new kinetic mass by $6/\theta$, closing the induction from zero
without assuming a mass-constrained projection.

The actual response satisfies
$\|\mQ\vx_k-\vb\|_2^2+\mu\|\vz_k-\vt\|_{\mQ}^2
\leq6\alpha\lambda_k$.
Consequently the squared response relative to $\vu_k$ is at most
$14\alpha\lambda_k$. Both the raw error and kinetic floor error are
downward; on a selected positive coordinate they can only reduce the
signed flow. Its weighted telescope is unchanged. Repeating the core and
Cauchy--Schwarz proof with $H_K\leq14S_K$ gives the constant $60$.

At the freeze, the objective part of the Euclidean energy is at most
$6\lambda$, and its mirror part is at most $7\lambda$, by the same
RPPR bias calculation as in \cref{lem:revisit-freeze-energy}. Hence
$E\leq13\lambda$. The frozen rounded energy is at most
$\chi^j13\lambda+\Gamma$, with $\Gamma\leq\tau/2$.
Since $\chi^{1/\theta}<1/2$, the stopping count makes the first term
at most $\tau/2$. The semantic clipping proof follows. The rounded
geometric sum supplies the work assertion.
\end{proof}

\paragraph{Rounded RPPR objective completion.}
\Cref{cor:revisit-rppr} also survives these rounded operations.
Use its $\rho,\delta$, put $\tau=\alpha\delta^2/2$, and choose $h$
by \cref{eq:revisit-rounded-grid} with this $\tau$.
Take $q$ minimal with $2^q\geq52\rho/\delta^2$ and use $N=q/\theta$
frozen steps. The actual freeze energy is still $13\alpha\rho$,
the accumulated error is at most $\tau/2$, and the first term contracts
below $\tau/2$. Thus the same order and objective-gap certificate holds.
Since $52\rho/\delta^2<416/\epsobj$, the bit and operation bounds
have only logarithmic dependence on $1/\epsobj$ as well.

\Cref{cor:revisit-positive-residual} also applies to the actual rounded
output: replace the frozen tolerance by $\alpha\delta^2/2$, use the
same grid error budget, and perform the final clipping exactly on the
represented densities. Choosing $\delta=\alpha\rho/2$ gives the positive
residual certificate with only additional parameter logarithms. This final
clipping need not be dyadic; it uses only fixed rational input denominators.

\paragraph{Represented arithmetic and rebase charges.}
Write $\theta=1/T$ and $h=1/H$, where $T,H$ are powers of two.
The stored $X,L,\sigma,Z$ are integers over $H$; kinetic numerators
are multiples of $T$. Boundary primal densities are at most three,
normalized primal records at most six, and kinetic densities at most
$6/\theta$. Neighbor records have only an additional degree factor.
The degree-weighted key $d_i k_i$ has a common denominator formed from
the fixed input denominator of $\alpha$, $H$, $T(T+1)$, and the current
scale numerator. An unweighted comparison uses only two individual
degrees; no least common multiple over degrees is formed. The threshold
also uses the fixed denominator of $g$ (and the initial $d_v$).
Floors immediately return every updated state to the grid.

Crossing the scale from one to below $1/2$ requires at least
$1/[2(\theta+h)]$ steps. There are therefore at most $4\theta K+1$
rebases, each a scalar/tree pass over retained records and no row scans.
The iteration count is $O(\theta^{-1}\log(1/\epsppr))$, and the union
of represented labels is charged by \cref{eq:revisit-rounded-work}.
Thus these passes add only a logarithmic factor.
For rational inputs the total local operation count remains
$\widetilde O(1/(\epsppr\sqrt\alpha))$; every integer has
\[
 O\bigl(L_{\rm in}+\log(1/(\alpha\epsppr))
                +\log(d_{\max}+2)\bigr)
\]
bits, including encountered label encodings in $L_{\rm in}$.
Schoolbook integer arithmetic multiplies the operation count by the square
of this bit bound. Exact threshold ties need no positive activation margin.
